%% 1-6 かずの 線：数直線（第11回で 説明を ふやした）
%% 第1巻のきまり：15 pt・行送り28 pt、分かち書き（文節の間に半角の空白）、漢字はすべてルビ、式は全角の＋－＝で書く

%% ---- この章だけで使う絵の部品（頭文字 \bkqf…） ----
\makeatletter
% 数直線  \bkqfline[数字の字]{(x,y)}{目もり1つの長さ}{目もりの数}{1目もりの大きさ}{数字を書く間（目もり いくつごと）}  % noruby
%   (x,y) は 0の目もり。右はしに やじるし。  % noruby
\newcommand\bkqfline[6][\small]{\begin{scope}[shift={#2}]
  \draw[bkInk, line width=1pt, -{Stealth[length=6pt, width=5pt]}] (-3,0) -- (#4*#3+7,0);
  \foreach \i in {0,...,#4} {\pgfmathtruncatemacro\bkqf@lab{mod(\i,#6)==0}
    \ifnum\bkqf@lab=1 \draw[bkInk, line width=0.9pt] (\i*#3,-1.9) -- (\i*#3,1.9);
      \pgfmathtruncatemacro\bkqf@v{\i*#5}
      \node[anchor=north, font=\sffamily\bfseries#1, text=bkInk, inner sep=0pt] at (\i*#3,-2.9) {\bkqf@v};
    \else\draw[bkInk, line width=0.6pt] (\i*#3,-1.3) -- (\i*#3,1.3);\fi}
\end{scope}}
% もどる やじるし（左へ）  \bkqfback{(x,y)}{目もり1つの長さ}{はじめ}{もどる数}{色}  % noruby
%   (x,y) は 0の目もりの真下で、弧の上はしの高さ。弧は下に ふくらむ。  % noruby
\newcommand\bkqfback[5]{\begin{scope}[shift={#1}]
  \pgfmathtruncatemacro\bkqf@e{#3-#4+1}
  \foreach \i in {#3,...,\bkqf@e} \draw[#5, line width=1.3pt, -{Stealth[length=5pt]}] (\i*#2-0.6,0) .. controls (\i*#2-0.25*#2,-6) and (\i*#2-0.75*#2,-6) .. (\i*#2-#2+0.6,-0.2);
\end{scope}}
% 数直線の 上の しるし（上から さす やじるしと 文字）  \bkqfmark{(x,y)}{文字}{色}   (x,y) は 線の上の点  % noruby
\newcommand\bkqfmark[3]{\begin{scope}[shift={#1}]
  \draw[#3, line width=1.4pt, -{Stealth[length=4.5pt]}] (0,6.6) -- (0,2.4);
  \node[circle, fill=#3, text=white, font=\sffamily\bfseries\small, minimum size=5.4mm, inner sep=0pt] at (0,9.6) {#2};
\end{scope}}
% かえる  \bkqffrog{(x,y)}{はば}   (x,y) は 足もとの まん中、右を むく  % noruby
\newcommand\bkqffrog[2]{\begin{scope}[shift={#1}, scale=#2/12]
  \fill[bkLeaf!85!black] (-4.4,0.6) ellipse (2 and 0.9) (3.8,0.6) ellipse (2 and 0.9);
  \fill[bkLeaf] (0,4) ellipse (5 and 3.8);
  \fill[bkYellow!35] (0.6,3) ellipse (3 and 2.2);
  \fill[bkLeaf] (-1.6,7.4) circle (1.8) (2.4,7.4) circle (1.8);
  \fill[white] (-1.6,7.6) circle (1.2) (2.4,7.6) circle (1.2);
  \fill[bkInk] (-1.2,7.6) circle (0.55) (2.8,7.6) circle (0.55);
  \fill[bkRed!45] (-3.6,4.6) circle (0.7) (4,4.6) circle (0.7);
  \draw[bkInk, line width=0.5pt] (-0.4,5.2) .. controls (0.4,4.4) and (1.4,4.4) .. (2,5.2);
\end{scope}}
% すごろくの カード  \bkqfcard{(x,y)}{文字}{色}  % noruby
\newcommand\bkqfcard[3]{\node[draw=#3, line width=1pt, fill=white, rounded corners=1.5mm, text=#3!80!black,
  font=\sffamily\bfseries\small, inner xsep=2mm, inner ysep=1.2mm] at #1 {#2};}
% とんだ 回数の 番号  \bkqfnums{(x,y)}{目もり1つの長さ}{はじめ}{いくつ}{色}{1:右へ / -1:左へ}  % noruby
%   (x,y) は 0の目もりの真上（または真下）で、番号を ならべる 高さ。  % noruby
\newcommand\bkqfnums[6]{\begin{scope}[shift={#1}]
  \foreach \k in {1,...,#4} {\node[circle, fill=white, draw=#5, line width=0.6pt, font=\sffamily\bfseries\scriptsize, inner sep=0.5pt, minimum size=3.6mm, text=#5] at ({(#3+#6*(\k-0.5))*#2},0) {\k};}
\end{scope}}
\makeatother

\bkchapter{1}{6}{かずの \R{線}{せん}：\R{数直線}{すう|ちょく|せん}}{\bklevel{1}}{かずの \R{線}{せん}：\R{数直線}{すう|ちょく|せん}}
\bkmichishirube{1-2,1-3,1-4,1-5}{1-7,1-8,1-11,2-2}

\begin{goal}
\begin{itemize}[label={\color{bkTheme}●}, leftmargin=1.4\zw, itemsep=0pt, topsep=0pt]
\item \R{数直線}{すう|ちょく|せん}の \R{上}{うえ}で、 たし\R{算}{ざん}と ひき\R{算}{ざん}が できる。
\item 2つの かずの ちがいを、 \R{数直線}{すう|ちょく|せん}で みつけられる。
\item 1\R{目}{め}もりが 10や 100の \R{数直線}{すう|ちょく|せん}を よめる。
\end{itemize}
\end{goal}

\begin{hatena}
\begin{bkfigure}
\begin{tikzpicture}[x=1mm, y=1mm]
  \fill[white, rounded corners=3mm] (-7,-31) rectangle (120,25);
  \bkruler{(0,0)}{9}{12}
  \bkhop{(0,0)}{9}{3}{5}{bkBlue}
  \bkqfback{(0,-10.5)}{9}{8}{2}{bkRed}
  \bkqffrog{(26,10.5)}{10}
  \draw[bkLeaf, line width=0.8pt, line cap=round] (20.5,11) -- (19,10) (20,13.5) -- (18.3,13.5);
  \bkqfcard{(58,18)}{5つ すすむ}{bkBlue}
  \bkqfcard{(78,-23)}{2つ もどる}{bkRed}
  \node[font=\sffamily\bfseries\Large, text=bkTheme] at (54,-23) {？};
\end{tikzpicture}
\end{bkfigure}
\noindent かえるが、 ものさしの 3から 5つ すすみました。\\
そこから 2つ もどると、 どこに つくかな？\\
もどる ときは、 どんな しきに なるのかな？
\end{hatena}

\clearpage
%% ---------------- p.67 やってみよう ----------------
\begin{yattemiyou}[かずの \R{線}{せん}で すごろく]
\begin{bkfigure}
\begin{tikzpicture}[x=1mm, y=1mm]
  \bkqfline{(4,0)}{11}{10}{1}{1}
  \bkohajiki{(4,4.5)}{2.8}{bkAppleRed}
  \node[anchor=south, font=\sffamily\bfseries\footnotesize, text=bkAppleRed!80!black] at (4,7.8) {スタート};
  \bkqfcard{(18,-15)}{4 すすむ}{bkBlue}
  \bkqfcard{(46,-15)}{1 もどる}{bkRed}
  \bkqfcard{(74,-15)}{6 すすむ}{bkBlue}
  \bkqfcard{(102,-15)}{2 もどる}{bkRed}
  \foreach \x/\n in {18/1,46/2,74/3,102/4} {\node[anchor=east, font=\sffamily\bfseries\footnotesize, text=bkGray] at (\x-11,-15) {\n};}
\end{tikzpicture}
\end{bkfigure}
\begin{enumerate}[label={\bkcircled[bkGreen]{\arabic*}}, leftmargin=2.2\zw, itemsep=2pt, topsep=2pt]
\item おはじきを 0に おこう。
\item カードの とおりに、 1まいずつ すすんだり もどったり しよう。 とまった ところの かずを かこう。 わからなく なったら、 0から やりなおして いいよ。
\end{enumerate}
\par\vspace{2pt}
\noindent\hspace*{1\zw}\begin{tblr}{colspec={*{2}{Q[l,m,wd=60mm]}}, colsep=0pt, rowsep=3pt}
1まいめの あと\enspace\bkkakikomi & 2まいめの あと\enspace\bkkakikomi \\
3まいめの あと\enspace\bkkakikomi & 4まいめの あと\enspace\bkkakikomi \\
\end{tblr}
\begin{enumerate}[label={\bkcircled[bkGreen]{\arabic*}}, leftmargin=2.2\zw, itemsep=2pt, topsep=4pt, start=3]
\item ろうかに テープで かずの \R{線}{せん}を かいて、 じぶんが こまに なって あるいても できるよ。
\end{enumerate}
\begin{bkfigure}
\begin{tikzpicture}[x=1mm, y=1mm]
  \fill[bkGray!18, rounded corners=2mm] (-4,-8) rectangle (104,22);
  \draw[bkYellow!80!black, line width=2.4pt] (0,0) -- (96,0);
  \foreach \i in {0,...,8} {\draw[bkYellow!80!black, line width=1.6pt] (\i*12,-2.5) -- (\i*12,2.5);
    \node[anchor=north, font=\sffamily\bfseries\footnotesize, text=bkInk] at (\i*12,-3.2) {\i};}
  \bkkid{(36,1.5)}{17}{bkGreen}{bkHair}
  \draw[bkBlue, line width=1.3pt, -{Stealth[length=5pt]}] (40,13) .. controls (43,18) and (47,17) .. (48,6);
  \node[anchor=west, font=\sffamily\bfseries\footnotesize, text=bkBlue!80!black] at (51,14) {1ぽ すすむ};
\end{tikzpicture}
\end{bkfigure}
\end{yattemiyou}

\bkfukidashi{delta}{すすむと かずが \R{大}{おお}きく なって、 もどると \R{小}{ちい}さく なったね。}

\clearpage
%% ---------------- 数学の道具（第11回：説明の段を ふやした） ----------------
\bksection{doguu}{\R{数学}{すう|がく}の \R{道具}{どう|ぐ}}

\bksubsection{1}{ものさしは かずの \R{線}{せん}}

\label{bk:1-6:sen}ものさしから、 \R{目}{め}もりと かずだけを とりだして みよう。

\begin{bkfigure}
\begin{tikzpicture}[x=1mm, y=1mm]
  \bkruler{(0,16)}{11}{10}
  \foreach \i in {0,...,10} {\draw[bkGray!70, line width=0.5pt, dash pattern=on 1.2pt off 1.2pt] (\i*11,6.5) -- (\i*11,2.2);}
  \bkqfline{(0,0)}{11}{10}{1}{1}
  \node[anchor=west, font=\sffamily\bfseries\footnotesize, text=bkInk] at (119,0) {つづく};
\end{tikzpicture}
\end{bkfigure}
\noindent まっすぐな \R{線}{せん}に、 おなじ あいだで \R{目}{め}もりと かずを つけた ものを、 \R{数直線}{すう|ちょく|せん}と いいます。\bkidx{すうちょくせん}{\R{数直線}{すう|ちょく|せん}}\bkidx{かずのせん}{かずの \R{線}{せん}}\bkidx{めもり}{\R{目}{め}もり}\\
\R{数直線}{すう|ちょく|せん}の きまりは 3つ です。
\begin{enumerate}[label={\bkcircled[bkTheme]{\arabic*}}, leftmargin=2.2\zw, itemsep=0pt, topsep=2pt]
\item \R{左}{ひだり}はしは 0から はじまる（1-3の ものさしと おなじ）。
\item \R{目}{め}もりと \R{目}{め}もりの あいだは、 どこも おなじ ながさ。
\item \R{右}{みぎ}へ 1つ いくと、 かずが 1 ふえる。
\end{enumerate}
\begin{bkfigure}
\begin{tikzpicture}[x=1mm, y=1mm]
  \bkqfline{(0,0)}{11}{10}{1}{1}
  \foreach \k in {1,...,10} {\draw[bkBlue, line width=1.1pt, -{Stealth[length=4pt]}] ({(\k-1)*11+0.6},2.4) .. controls ({(\k-0.8)*11},8) and ({(\k-0.2)*11},8) .. ({\k*11-0.6},2.6);
    \node[font=\sffamily\bfseries\scriptsize, text=bkBlue!85!black] at ({(\k-0.5)*11},10.3) {＋1};}
  \foreach \k in {2,7} {\draw[bkGreen!70!black, line width=0.7pt, decorate, decoration={brace, amplitude=1.4mm, mirror}] (\k*11+0.5,-8.5) -- (\k*11+10.5,-8.5);}
  \node[font=\sffamily\bfseries\scriptsize, text=bkGreen!60!black] at (27.5,-14) {おなじ};
  \node[font=\sffamily\bfseries\scriptsize, text=bkGreen!60!black] at (82.5,-14) {おなじ};
  \draw[bkTheme, line width=1pt, -{Stealth[length=4pt]}] (-6,-13) -- (-1,-7);
  \node[anchor=north, font=\sffamily\bfseries\scriptsize, text=bkTheme] at (-6,-13.5) {0から};
\end{tikzpicture}
\end{bkfigure}
\R{右}{みぎ}はしの やじるしは、 「まだ \R{右}{みぎ}へ つづくよ」の しるしです。

\clearpage
\noindent{\sffamily\bfseries\color{bkTheme}なぜ 「おなじ あいだ」に するの？}\par
\R{目}{め}もりの あいだが ばらばらだったら、 どう なるかな？

\begin{bkfigure}
\begin{tikzpicture}[x=1mm, y=1mm]
  % × ばらばらの 目もり
  \node at (-6,0) {\bkbatsu};
  \draw[bkInk, line width=1pt, -{Stealth[length=6pt, width=5pt]}] (0,0) -- (122,0);
  \foreach \x/\v in {4/0,10/1,40/2,48/3,80/4,88/5,112/6} {\draw[bkInk, line width=0.9pt] (\x,-1.9) -- (\x,1.9);
    \node[anchor=north, font=\sffamily\bfseries\small, text=bkInk, inner sep=0pt] at (\x,-2.9) {\v};}
  \draw[bkBlue, line width=1.3pt, -{Stealth[length=5pt]}] (10.6,2.4) .. controls (18,12) and (32,12) .. (39.4,2.6);
  \draw[bkBlue, line width=1.3pt, -{Stealth[length=5pt]}] (40.6,2.4) .. controls (42.5,8) and (45.5,8) .. (47.4,2.6);
  \node[anchor=south, font=\sffamily\bfseries\footnotesize, text=bkBlue!80!black] at (25,10.5) {1 すすむ};
  \node[anchor=south, font=\sffamily\bfseries\footnotesize, text=bkBlue!80!black] at (52,8.5) {1 すすむ};
  % ○ おなじ あいだ
  \begin{scope}[shift={(0,-26)}]
  \node at (-6,0) {\bkmaru};
  \bkqfline{(4,0)}{18}{6}{1}{1}
  \bkhop{(4,-3)}{18}{1}{2}{bkBlue}
  \node[anchor=west, font=\sffamily\bfseries\footnotesize, text=bkBlue!80!black] at (64,7) {どちらも おなじ ながさ};
  \end{scope}
\end{tikzpicture}
\end{bkfigure}
\noindent \bkbatsu の \R{線}{せん}では、 おなじ 「1 すすむ」 なのに、 ながい ときと みじかい ときが あります。 これでは、 ながさを \R{見}{み}て かずを つかむ ことが できません。\\
\bkmaru の \R{線}{せん}は、 「1 すすむ」が いつも おなじ ながさ。 だから、 \R{目}{め}もりを かぞえれば、 いくつ すすんだかが わかります。

\par\vspace{6pt}
\noindent{\sffamily\bfseries\color{bkTheme}\R{右}{みぎ}に ある ほど \R{大}{おお}きい}\par
\R{右}{みぎ}へ 1つ いくと 1 ふえるので、 \R{右}{みぎ}に ある かずの ほうが \R{大}{おお}きく なります。

\begin{bkfigure}
\begin{tikzpicture}[x=1mm, y=1mm]
  \bkqfline{(0,0)}{11}{10}{1}{1}
  \bkqfmark{(44,0)}{4}{bkBlue}
  \bkqfmark{(77,0)}{7}{bkTheme}
  \draw[bkGray, line width=0.9pt, -{Stealth[length=4pt]}] (50,9.6) -- (71,9.6);
  \node[anchor=south, font=\sffamily\bfseries\footnotesize, text=bkGray] at (60.5,10.8) {\R{右}{みぎ}に ある};
\end{tikzpicture}
\end{bkfigure}
\noindent 7は 4より \R{右}{みぎ}に あるので、 7の ほうが \R{大}{おお}きい かずです。

\clearpage
\bksubsection{2}{たし\R{算}{ざん}は \R{右}{みぎ}へ すすむ}

\label{bk:1-6:susumu}\bkhatenaref の かえるは、 3から 5つ すすみました。 \R{数直線}{すう|ちょく|せん}で \R{見}{み}て みよう。

\begin{bkfigure}
\begin{tikzpicture}[x=1mm, y=1mm]
  \bkqfline{(0,0)}{11}{10}{1}{1}
  \bkhop{(0,-3)}{11}{3}{5}{bkBlue}
  \bkqfnums{(0,9.6)}{11}{3}{5}{bkBlue}{1}
  \node[anchor=west, font=\sffamily\bfseries\small, text=bkBlue!80!black] at (92,6) {5つ すすむ};
\end{tikzpicture}
\end{bkfigure}
\begin{enumerate}[label={\bkcircled[bkTheme]{\arabic*}}, leftmargin=2.2\zw, itemsep=0pt, topsep=0pt]
\item 3から はじめる。 3は かぞえない。
\item 1かい とぶ ごとに 1 ふえる。 4、 5、 6、 7、 8。
\item 5かい とんで、 8に ついた。
\end{enumerate}
「5つ すすむ」は 「5 ふえる」 ことです。 ふえるのは たし\R{算}{ざん}でした（1-4）。
\bkbigshiki{3 ＋ 5 ＝ 8}{\R{右}{みぎ}へ 5 すすむ}

\noindent{\sffamily\bfseries\color{bkTheme}テープで \R{見}{み}ると}\par
0から 3までの テープと、 3から 8までの テープを ならべて みよう。

\begin{bkfigure}
\begin{tikzpicture}[x=1mm, y=1mm]
  \bkqfline{(0,0)}{11}{10}{1}{1}
  \bktape{(0,4)}{33}{bkBlue}{3}
  \bktape{(33,4)}{55}{bkOrange}{5}
  \bkbrace{(0,12)}{(88,12)}{8}{1}
\end{tikzpicture}
\end{bkfigure}
\noindent 3の ながさと 5の ながさを つなぐと、 8の ながさに なります。 たし\R{算}{ざん}は、 「ながさを つなぐ」 ことでも あるのです。

\clearpage
\bksubsection{3}{ひき\R{算}{ざん}は \R{左}{ひだり}へ もどる}

\label{bk:1-6:modoru}かえるは、 8から 2つ もどりました。

\begin{bkfigure}
\begin{tikzpicture}[x=1mm, y=1mm]
  \bkqfline{(0,0)}{11}{10}{1}{1}
  \fill[bkBlue] (88,0) circle (1.6);
  \bkqfback{(0,-9.5)}{11}{8}{2}{bkRed}
  \bkqfnums{(0,-19.5)}{11}{8}{2}{bkRed}{-1}
  \node[anchor=west, font=\sffamily\bfseries\small, text=bkRed!80!black] at (92,-13) {2つ もどる};
\end{tikzpicture}
\end{bkfigure}
\begin{enumerate}[label={\bkcircled[bkTheme]{\arabic*}}, leftmargin=2.2\zw, itemsep=0pt, topsep=0pt]
\item 8から はじめる。
\item 1かい もどる ごとに 1 へる。 7、 6。
\item 2かい もどって、 6に ついた。
\end{enumerate}
「2つ もどる」は 「2 へる」 ことです。 へるのは ひき\R{算}{ざん}でした（1-5）。
\bkbigshiki{8 － 2 ＝ 6}{\R{左}{ひだり}へ 2 もどる}
\bkhatenaref の こたえは 6。 もどる ときの しきは、 ひき\R{算}{ざん}です。

\par\vspace{4pt}
\noindent{\sffamily\bfseries\color{bkTheme}もどって たしかめよう}\par
6から、 こんどは 2つ すすんで みよう。

\begin{bkfigure}
\begin{tikzpicture}[x=1mm, y=1mm]
  \bkqfline{(0,0)}{11}{10}{1}{1}
  \bkhop{(0,-3)}{11}{6}{2}{bkBlue}
  \node[anchor=west, font=\sffamily\bfseries\small, text=bkBlue!80!black] at (92,6) {2つ すすむ};
\end{tikzpicture}
\end{bkfigure}
\noindent 6 ＋ 2 ＝ 8で、 もとの 8に もどりました。 2つ もどって、 2つ すすむと、 はじめの ところに かえります。 ひき\R{算}{ざん}の こたえは、 たし\R{算}{ざん}で たしかめられます。

\clearpage
\bksubsection{4}{2つの かずの ちがい}

\label{bk:1-6:aida}ソラは 3、 ミナは 8の ところに います。 ふたりは いくつ はなれて いる？

\begin{bkfigure}
\begin{tikzpicture}[x=1mm, y=1mm]
  \bkqfline{(0,0)}{11}{10}{1}{1}
  \bkChar{sora}{(33,1.4)}{16mm}
  \bkChar{mina}{(88,1.4)}{16mm}
  \node[font=\sffamily\bfseries\footnotesize, text=bkRoute2!80!black] at (33,19.5) {ソラ};
  \node[font=\sffamily\bfseries\footnotesize, text=bkAppleRed!80!black] at (88,19.5) {ミナ};
  \foreach \i in {3,...,7} {\fill[bkOrange!70] (\i*11+1,-0.9) rectangle (\i*11+10,0.9);}
  \bkbrace{(33,-9)}{(88,-9)}{5}{-1}
\end{tikzpicture}
\end{bkfigure}
\noindent{\sffamily\bfseries しかた1：かぞえる}\quad 3から 8まで、 1つずつ すすむと 5かい。 オレンジの あいだも 5つ。 ふたりは \mbox{5 はなれて} います。

\par\vspace{4pt}
\noindent{\sffamily\bfseries しかた2：0からの テープで くらべる}\par
ミナの 8は 「0から 8つ」、 ソラの 3は 「0から 3つ」です。

\begin{bkfigure}
\begin{tikzpicture}[x=1mm, y=1mm]
  \bkqfline{(0,0)}{11}{10}{1}{1}
  \bktape{(0,13)}{88}{bkAppleRed}{ミナ 8}
  \bktape{(0,4)}{33}{bkRoute2}{ソラ 3}
  \fill[bkOrange!35] (33,4) rectangle (88,11);
  \draw[bkOrange!80!black, line width=0.8pt, dash pattern=on 2pt off 1.5pt] (33,4) rectangle (88,11);
  \node[font=\sffamily\bfseries\small, text=bkOrange!70!black] at (60.5,7.5) {ちがい};
\end{tikzpicture}
\end{bkfigure}
\noindent 8の テープから 3の テープを とると、 のこりが ちがいです。 ちがいは ひき\R{算}{ざん}でした（1-5）。
\bkbigshiki{8 － 3 ＝ 5}{8と 3の ちがいは 5}
\noindent かぞえても、 ひき\R{算}{ざん}でも、 おなじ 5に なりました。 かずが \R{大}{おお}きく なると かぞえるのは たいへん。 そんな ときは ひき\R{算}{ざん}が べんりです。

\clearpage
\bksubsection{5}{1\R{目}{め}もりが 10の \R{数直線}{すう|ちょく|せん}}

\label{bk:1-6:ookii}100の ような \R{大}{おお}きな かずを、 1ずつの \R{目}{め}もりで かくと、 \R{線}{せん}が とても ながく なって \R{本}{ほん}に はいりません。\\
そこで、 10の まとまり（1-2）で \R{目}{め}もりを つけます。

\begin{bkfigure}
\begin{tikzpicture}[x=1mm, y=1mm]
  \bkqfline[\footnotesize]{(0,0)}{11}{10}{10}{1}
  \foreach \i in {0,...,6} {\foreach \j in {0,...,9} {\filldraw[fill=bkRoute4!70, draw=bkRoute4!50!black, line width=0.3pt] (\i*11+\j*1.1,3.2) rectangle ++(1.1,1.6);}}
  \bkbrace{(0,6)}{(11,6)}{1が 10こ}{1}
  \bkqfmark{(77,0)}{A}{bkTheme}
\end{tikzpicture}
\end{bkfigure}
\noindent 1つの あいだに、 1が 10こ はいって います。 だから、 \R{右}{みぎ}へ 1\R{目}{め}もり すすむと 10 ふえます。 この 「\R{目}{め}もり 1つ\R{分}{ぶん}の \R{大}{おお}きさ」を、 1\R{目}{め}もりと いいます。\bkidx{ひとめもり}{1\R{目}{め}もり}\\
A は、 0から 7\R{目}{め}もり。 10が 7こで 70です（1-2）。

\par\vspace{2pt}
\noindent{\sffamily\bfseries\color{bkTheme}1\R{目}{め}もりが 100の とき}\par
1\R{目}{め}もりが 100なら、 0から 7\R{目}{め}もりの ところは 100が 7こで 700です。

\begin{bkfigure}
\begin{tikzpicture}[x=1mm, y=1mm]
  \bkqfline[\footnotesize]{(0,0)}{11}{10}{100}{1}
  \bkqfmark{(77,0)}{A}{bkTheme}
  \foreach \k in {1,...,7} {\draw[bkBlue, line width=1.1pt, -{Stealth[length=4pt]}] ({(\k-1)*11+0.6},-9) .. controls ({(\k-0.8)*11},-14) and ({(\k-0.2)*11},-14) .. ({\k*11-0.6},-9.2);
    \node[font=\sffamily\bfseries\scriptsize, text=bkBlue!85!black] at ({(\k-0.5)*11},-16.4) {\k};}
  \node[anchor=west, font=\sffamily\bfseries\footnotesize, text=bkBlue!80!black] at (80,-13) {100が 7こ};
\end{tikzpicture}
\end{bkfigure}
\noindent おなじ 7\R{目}{め}もりでも、 1\R{目}{め}もりが 10なら 70、 100なら 700。 1\R{目}{め}もりの \R{大}{おお}きさで、 よみかたが かわります。

\clearpage

\bksubsection{6}{かずが とびとびに かいて ある とき}

\label{bk:1-6:tobitobi}\R{目}{め}もりが こまかいと、 ぜんぶの \R{目}{め}もりに かずは かけません。 0と 50と 100だけ かいて ある \R{数直線}{すう|ちょく|せん}を \R{見}{み}て みよう。

\begin{bkfigure}
\begin{tikzpicture}[x=1mm, y=1mm]
  \bkqfline[\footnotesize]{(0,0)}{11}{10}{10}{5}
  \foreach \k in {1,...,5} {\draw[bkBlue, line width=1.2pt, -{Stealth[length=4.5pt]}] ({(\k-1)*11+0.6},2.4) .. controls ({(\k-0.8)*11},8) and ({(\k-0.2)*11},8) .. ({\k*11-0.6},2.6);
    \pgfmathtruncatemacro\bkqfv{10*\k}\node[font=\sffamily\bfseries\scriptsize, text=bkBlue!85!black] at ({(\k-0.5)*11},10.4) {\bkqfv};}
\end{tikzpicture}
\end{bkfigure}
\begin{enumerate}[label={\bkcircled[bkTheme]{\arabic*}}, leftmargin=2.2\zw, itemsep=0pt, topsep=0pt]
\item かずの ある 0と 50の あいだを かぞえる。 \mbox{5つ ある。}
\item 1ずつ ふやすと、 1、 2、 3、 4、 5で、 50に ならない。 \bkbatsu
\item 10ずつ ふやすと、 10、 20、 30、 40、 50で、 50に ぴったり。 \bkmaru
\end{enumerate}
だから、 この \R{数直線}{すう|ちょく|せん}の 1\R{目}{め}もりは 10です。 （1\R{目}{め}もりの みつけかたは、 1-11で もっと くわしく やります。）

\begin{taisetsu}
\R{右}{みぎ}へ すすむのは たし\R{算}{ざん}、 \R{左}{ひだり}へ もどるのは ひき\R{算}{ざん}。\\
2つの かずの ちがいは、 \R{大}{おお}きい かず － \R{小}{ちい}さい かず。\\
よむ ときは、 まず 1\R{目}{め}もりの \R{大}{おお}きさを \R{見}{み}よう。
\end{taisetsu}

\begin{column}{\R{数直線}{すう|ちょく|せん}は どこまで つづく？}
\R{右}{みぎ}へは、 100、 1000、 10000… どこまでも つづきます。 では、 0より \R{左}{ひだり}は？ さむい \R{日}{ひ}の おんどけいには、 0より \R{下}{した}の \R{目}{め}もりが あります。 0より \R{小}{ちい}さい かずは、 この \R{本}{ほん}の 2\R{巻}{かん}で まなびます。
\end{column}

\clearpage
%% ---------------- 例題 ----------------

\begin{reidai}[\R{右}{みぎ}へ すすむ]
8から 5 すすむと、 どこ？
\begin{bkfigure}
\begin{tikzpicture}[x=1mm, y=1mm]
  \bkqfline[\footnotesize]{(0,0)}{7.2}{15}{1}{1}
  \bkhop{(0,-3)}{7.2}{8}{5}{bkBlue}
  \bkqfnums{(0,9.4)}{7.2}{8}{5}{bkBlue}{1}
\end{tikzpicture}
\end{bkfigure}
\kangaekata \R{右}{みぎ}へ すすむから、 たし\R{算}{ざん}。 9、 10、 11、 12、 13と 5かい とぶ。 10を こえても おなじ。

\shiki 8 ＋ 5 ＝ 13 \kotaenext 13の ところ
\end{reidai}

\begin{jibun}
\R{数直線}{すう|ちょく|せん}の 6から 5 すすむと、 どこ？
\bkshikiran[9\zw]
\bkjans{6 ＋ 5 ＝ 11\quad 11の ところ}
\end{jibun}

\begin{reidai}[\R{左}{ひだり}へ もどる]
12から 4 もどると、 どこ？
\begin{bkfigure}
\begin{tikzpicture}[x=1mm, y=1mm]
  \bkqfline[\footnotesize]{(0,0)}{7.2}{15}{1}{1}
  \bkqfback{(0,-8)}{7.2}{12}{4}{bkRed}
  \bkqfnums{(0,-16.5)}{7.2}{12}{4}{bkRed}{-1}
  \fill[bkRed] (86.4,0) circle (1.6);
\end{tikzpicture}
\end{bkfigure}
\kangaekata \R{左}{ひだり}へ もどるから、 ひき\R{算}{ざん}。 11、 10、 9、 8と 4かい もどる。

\shiki 12 － 4 ＝ 8 \kotaenext 8の ところ

\end{reidai}

\begin{jibun}
\R{数直線}{すう|ちょく|せん}の 13から 6 もどると、 どこ？
\bkshikiran[9\zw]
\bkjans{13 － 6 ＝ 7\quad 7の ところ}
\end{jibun}

\begin{reidai}[ちがいを みつける]
4と 11は、 いくつ ちがう？
\begin{bkfigure}
\begin{tikzpicture}[x=1mm, y=1mm]
  \bkqfline[\footnotesize]{(0,0)}{7.2}{15}{1}{1}
  \fill[bkTheme] (28.8,0) circle (1.6) (79.2,0) circle (1.6);
  \foreach \i in {4,...,10} {\fill[bkOrange!70] (\i*7.2+0.8,-0.9) rectangle (\i*7.2+6.4,0.9);}
  \bkbrace{(28.8,2.4)}{(79.2,2.4)}{7}{1}
\end{tikzpicture}
\end{bkfigure}
\kangaekata 4から 11まで、 オレンジの あいだを かぞえると 7つ。 ちがいは ひき\R{算}{ざん}で、 \R{大}{おお}きい ほうから \R{小}{ちい}さい ほうを ひく。

\shiki 11 － 4 ＝ 7 \kotaenext 7
\end{reidai}

\begin{jibun}
5と 13は、 いくつ ちがう？
\bkshikiran
\bkjans{13 － 5 ＝ 8\quad 8}
\end{jibun}

\begin{reidai}[1\R{目}{め}もりが 10]
A と B の かずは いくつ？
\begin{bkfigure}
\begin{tikzpicture}[x=1mm, y=1mm]
  \bkqfline[\footnotesize]{(0,0)}{10.5}{10}{10}{5}
  \bkqfmark{(31.5,0)}{A}{bkBlue}
  \bkqfmark{(84,0)}{B}{bkBlue}
  \foreach \k/\v in {1/10,2/20,3/30} {\draw[bkBlue, line width=1.1pt, -{Stealth[length=4pt]}] ({(\k-1)*10.5+0.6},-9) .. controls ({(\k-0.8)*10.5},-14) and ({(\k-0.2)*10.5},-14) .. ({\k*10.5-0.6},-9.2);
    \node[font=\sffamily\bfseries\scriptsize, text=bkBlue!85!black] at ({(\k-0.5)*10.5},-16.4) {\v};}
  \foreach \k/\v in {6/60,7/70,8/80} {\draw[bkBlue, line width=1.1pt, -{Stealth[length=4pt]}] ({(\k-1)*10.5+0.6},-9) .. controls ({(\k-0.8)*10.5},-14) and ({(\k-0.2)*10.5},-14) .. ({\k*10.5-0.6},-9.2);
    \node[font=\sffamily\bfseries\scriptsize, text=bkBlue!85!black] at ({(\k-0.5)*10.5},-16.4) {\v};}
\end{tikzpicture}
\end{bkfigure}
\kangaekata 0と 50の あいだは 5つ。 10ずつで 50に ぴったりだから、 1\R{目}{め}もりは 10。\\
A は 0から 3\R{目}{め}もり。 10、 20、 30。\\
B は 50から 3\R{目}{め}もり。 60、 70、 80。

\kotae A\enspace 30\qquad B\enspace 80
\end{reidai}

\begin{jibun}
C と D の かずは？\hfill C\enspace\bkkakikomi[22mm]\hspace{1.5\zw}D\enspace\bkkakikomi[22mm]
\begin{bkfigure}
\begin{tikzpicture}[x=1mm, y=1mm]
  \bkqfline[\footnotesize]{(0,0)}{10.5}{10}{10}{5}
  \bkqfmark{(42,0)}{C}{bkOrange}
  \bkqfmark{(94.5,0)}{D}{bkOrange}
\end{tikzpicture}
\end{bkfigure}
\noindent{\small\color{bkOrange!58!black}ヒント：}0と 50の あいだは\enspace\bkkakikomi[9mm]\enspace つ。 1\R{目}{め}もりは\enspace\bkkakikomi[12mm]
\bkjans{C\enspace 40\quad D\enspace 90（あいだは 5つ、 1\R{目}{め}もりは 10）}
\end{jibun}

\begin{reidai}[1\R{目}{め}もりが 100]
A と B の かずは いくつ？
\begin{bkfigure}
\begin{tikzpicture}[x=1mm, y=1mm]
  \bkqfline[\footnotesize]{(0,0)}{10.5}{10}{100}{5}
  \bkqfmark{(21,0)}{A}{bkBlue}
  \bkqfmark{(63,0)}{B}{bkBlue}
  \foreach \k/\v in {1/100,2/200,6/600} {\draw[bkBlue, line width=1.1pt, -{Stealth[length=4pt]}] ({(\k-1)*10.5+0.6},-9) .. controls ({(\k-0.8)*10.5},-14) and ({(\k-0.2)*10.5},-14) .. ({\k*10.5-0.6},-9.2);
    \node[font=\sffamily\bfseries\scriptsize, text=bkBlue!85!black] at ({(\k-0.5)*10.5},-16.4) {\v};}
\end{tikzpicture}
\end{bkfigure}
\kangaekata 0と 500の あいだは 5つ。 100ずつ ふやすと、 100、 200、 300、 400、 500で ぴったり。 だから 1\R{目}{め}もりは 100。\\
A は 0から 2\R{目}{め}もりで、 100が 2こ。\\
B は 500から 1\R{目}{め}もりで、 500と 100で 600。

\kotae A\enspace 200\qquad B\enspace 600
\end{reidai}

\begin{jibun}
A と B の かずは？\hfill A\enspace\bkkakikomi[22mm]\hspace{1.5\zw}B\enspace\bkkakikomi[22mm]
\begin{bkfigure}
\begin{tikzpicture}[x=1mm, y=1mm]
  \bkqfline[\footnotesize]{(0,0)}{10.5}{10}{100}{5}
  \bkqfmark{(42,0)}{A}{bkOrange}
  \bkqfmark{(73.5,0)}{B}{bkOrange}
\end{tikzpicture}
\end{bkfigure}
\noindent{\small\color{bkOrange!58!black}ヒント：}0と 500の あいだは\enspace\bkkakikomi[9mm]\enspace つ。 1\R{目}{め}もりは\enspace\bkkakikomi[12mm]
\bkjans{A\enspace 400\quad B\enspace 700（あいだは 5つ、 1\R{目}{め}もりは 100）}
\end{jibun}

\begin{reidai}[どこから すすんだ？]
かえるが ある ところから 4つ すすむと、 11に つきました。 どこから すすんだのかな？
\begin{bkfigure}
\begin{tikzpicture}[x=1mm, y=1mm]
  \bkqfline[\footnotesize]{(0,0)}{7.2}{15}{1}{1}
  \fill[bkBlue] (79.2,0) circle (1.6);
  \bkqfback{(0,-8)}{7.2}{11}{4}{bkRed}
  \bkqfnums{(0,-16.5)}{7.2}{11}{4}{bkRed}{-1}
  \node[circle, draw=bkTheme, line width=1pt, fill=white, text=bkTheme, font=\sffamily\bfseries\small, inner sep=0pt, minimum size=5mm] at (50.4,8) {？};
\end{tikzpicture}
\end{bkfigure}
\kangaekata すすむ \R{前}{まえ}の ところは、 11から 4つ もどれば わかる。 もどるのは ひき\R{算}{ざん}。

\shiki 11 － 4 ＝ 7 \kotaenext 7の ところ

\noindent{\sffamily\bfseries\small たしかめ}\par\noindent 7から 4つ すすむと、 ほんとうに 11に つくかな？
\begin{bkfigure}
\begin{tikzpicture}[x=1mm, y=1mm]
  \bkqfline[\footnotesize]{(0,0)}{7.2}{15}{1}{1}
  \bkhop{(0,-3)}{7.2}{7}{4}{bkBlue}
  \node[anchor=west, font=\sffamily\bfseries\small, text=bkBlue!80!black] at (84,6) {7 ＋ 4 ＝ 11\enspace\bkmaru};
\end{tikzpicture}
\end{bkfigure}
\end{reidai}

\begin{jibun}
ある ところから 5つ すすむと、 13に つきました。 どこから すすんだ？
\bkshikiran[9\zw]
\bkjans{13 － 5 ＝ 8\quad 8の ところ}
\end{jibun}

\clearpage
%% ---------------- p.72 まちがえやすいところ ----------------
\bksection{tsumazuki}{まちがえやすい ところ}

\begin{tsumazuki}{3から かぞえて しまう}
3から 4 すすむと？\hspace{2\zw}\bkbatsu\enspace 6\hspace{3\zw}\bkmaru\enspace 7
\begin{bkfigure}
\begin{tikzpicture}[x=1mm, y=1mm]
  \bkqfline[\footnotesize]{(0,0)}{11}{8}{1}{1}
  \bkhop{(0,-3)}{11}{3}{4}{bkBlue}
  \foreach \k in {1,...,4} {\node[circle, fill=white, draw=bkBlue, line width=0.6pt, font=\sffamily\bfseries\scriptsize, inner sep=0.6pt, text=bkBlue] at (\k*11+27.5,9) {\k};}
\end{tikzpicture}
\end{bkfigure}
\noindent 3の \R{目}{め}もりは かぞえない。 ぴょんと とんだ かずを かぞえよう。
\end{tsumazuki}

\begin{tsumazuki}{1\R{目}{め}もりは いつも 1？}
\noindent\begin{minipage}[c]{0.5\linewidth}
\begin{tikzpicture}[x=1mm, y=1mm]
  \bkqfline[\footnotesize]{(0,0)}{5.2}{10}{10}{5}
  \bkqfmark{(36.4,0)}{A}{bkTheme}
\end{tikzpicture}
\end{minipage}\hfill
\begin{minipage}[c]{0.46\linewidth}
A は？\hspace{1\zw}\bkbatsu\enspace 7\hspace{1.5\zw}\bkmaru\enspace 70\par
1\R{目}{め}もりが 10の ときも あるよ。
\end{minipage}
\end{tsumazuki}

\begin{tsumazuki}{ちがいで \R{目}{め}もりの \R{線}{せん}を かぞえて しまう}
3と 8の ちがいは？\hspace{2\zw}\bkbatsu\enspace 6\hspace{3\zw}\bkmaru\enspace 5
\begin{bkfigure}
\begin{tikzpicture}[x=1mm, y=1mm]
  \bkqfline[\footnotesize]{(0,0)}{9}{10}{1}{1}
  \foreach \i [count=\k] in {3,...,8} {\node[circle, fill=white, draw=bkRed, line width=0.6pt, font=\sffamily\bfseries\scriptsize, inner sep=0.5pt, minimum size=3.6mm, text=bkRed] at (\i*9,5.4) {\k};}
  \foreach \i [count=\k] in {3,...,7} {\fill[bkOrange!70] (\i*9+1,-0.9) rectangle (\i*9+8,0.9);
    \node[circle, fill=white, draw=bkGreen!70!black, line width=0.6pt, font=\sffamily\bfseries\scriptsize, inner sep=0.5pt, minimum size=3.6mm, text=bkGreen!60!black] at (\i*9+4.5,-11) {\k};}
  \node[anchor=west, font=\sffamily\bfseries\footnotesize, text=bkRed] at (99,5.4) {\R{線}{せん}は 6\R{本}{ほん}};
  \node[anchor=west, font=\sffamily\bfseries\footnotesize, text=bkGreen!60!black] at (99,-11) {あいだは 5つ};
\end{tikzpicture}
\end{bkfigure}
\noindent 3の \R{線}{せん}と 8の \R{線}{せん}を どちらも かぞえると、 1つ おおく なります。 \R{線}{せん}と \R{線}{せん}の 「あいだ」を かぞえよう。
\end{tsumazuki}



\clearpage
%% ---------------- p.73 つかわれているところ ----------------
\begin{butsuri}
\R{数直線}{すう|ちょく|せん}は、 まちの なかにも あります。
\begin{bkfigure}
\begin{tikzpicture}[x=1mm, y=1mm]
  % 電車の ろせん図  % noruby
  \fill[white, rounded corners=3mm] (0,0) rectangle (75,42);
  \draw[bkRoute6, line width=2.6pt, line cap=round] (6,22) -- (69,22);
  \foreach \k in {1,...,8} {\fill[white, draw=bkRoute6!80!black, line width=1.2pt] (\k*8-1,22) circle (2.2);
    \node[font=\sffamily\bfseries\footnotesize, text=bkInk] at (\k*8-1,16.5) {\k};}
  \fill[bkRoute6!85!black, rounded corners=1.2mm] (17,27) rectangle (29,33);
  \foreach \w in {19,23} {\fill[white] (\w,29.5) rectangle ++(2.6,2);}
  \draw[bk arrow, bkOrange, line width=1.3pt] (23,37) -- (55,37);
  \node[font=\sffamily\bfseries\small] at (37.5,6) {\R{電車}{でん|しゃ}の えきの ばんごう};
  % エレベーター  % noruby
  \begin{scope}[shift={(79,0)}]
  \fill[white, rounded corners=3mm] (0,0) rectangle (43,42);
  \foreach \f in {1,...,6} {\draw[bkGray!60, line width=0.6pt] (14,9+\f*5) -- (30,9+\f*5);
    \node[anchor=east, font=\sffamily\bfseries\scriptsize, text=bkInk] at (13,6.5+\f*5) {\f かい};}
  \draw[bkGray, line width=0.9pt] (14,9) rectangle (30,39);
  \fill[bkRoute2!80, rounded corners=0.8mm] (16,14.6) rectangle (24,18.4);
  \draw[bk arrow, bkOrange, line width=1.3pt] (27,16.5) -- (27,31.5);
  \node[font=\sffamily\bfseries\small] at (21.5,2.5) {エレベーター};
  \end{scope}
\end{tikzpicture}
\end{bkfigure}
\begin{itemize}[label={\color{bkPurple}●}, leftmargin=1.4\zw, itemsep=0pt, topsep=0pt]
\item えきの ばんごうも \R{数直線}{すう|ちょく|せん}。 3ばんから 7ばんへは、 7 － 3 ＝ 4で 4えき さきです。
\item エレベーターは、 たてに した \R{数直線}{すう|ちょく|せん}。 2かいから 5かいへは 3かい\R{分}{ぶん} のぼります。
\item すごろくや、 たいおんけいの \R{目}{め}もりも \R{数直線}{すう|ちょく|せん}です。
\end{itemize}
\end{butsuri}

\bkfukidashi{delta}{「どっちの むきに、 いくつ うごくか」を かずで あらわすと、 ロボットの うごきも かけるんだ。 「やじるしの \R{路線}{ろ|せん}」に のって いこう！}

\begin{bkfigure}
\begin{tikzpicture}[x=1mm, y=1mm]
  \draw[bkRoute6, line width=3pt] (10,0) -- (128,0);
  \node[anchor=west, inner sep=0pt] at (0,12) {\bkrouteicon[1.4em]{6}\hspace{0.3em}{\sffamily\bfseries\small\color{bkRoute6!70!black}やじるしの \R{路線}{ろ|せん}}};
  \foreach \x/\c/\t/\s in {10/bkVol1/{1-3}/{0},39.5/bkVol1/{1-6}/{\R{数直線}{すう|ちょく|せん}},69/bkVol2/{2-2}/{マイナス},98.5/bkVol3/{3-2}/{ざひょう},128/bkVol4/{4-10}/{ベクトル}} {
    \fill[white, draw=\c, line width=1.6pt] (\x,0) circle (2.6);
    \node[anchor=base, font=\sffamily\bfseries\footnotesize, text=\c!70!black] at (\x,-8.4) {\t};
    \node[anchor=base, font=\sffamily\footnotesize, text=bkInk] at (\x,-15.8) {\s};}
  \node[fill=bkTheme, text=white, rounded corners=1mm, font=\sffamily\bfseries\fontsize{8}{9}\selectfont, inner sep=1.5pt] at (39.5,5.5) {いま ここ};
\end{tikzpicture}
\end{bkfigure}

\clearpage
%% ---------------- p.74〜75 もんだい ----------------
\bkendproblems

\bkproblemgroup{きほん}{1}

\bkproblem \R{数直線}{すう|ちょく|せん}を つかって、 けいさんしよう。
\begin{bkfigure}
\begin{tikzpicture}[x=1mm, y=1mm]
  \bkqfline[\footnotesize]{(0,0)}{7.6}{15}{1}{1}
\end{tikzpicture}
\end{bkfigure}
\noindent\hspace*{1\zw}\begin{tblr}{colspec={*{2}{Q[l,m,wd=62mm]}}, colsep=0pt, rowsep=3pt}
（1）5 ＋ 4 ＝ \bkkakikomi & （2）7 ＋ 6 ＝ \bkkakikomi \\
（3）9 － 3 ＝ \bkkakikomi & （4）15 － 8 ＝ \bkkakikomi \\
\end{tblr}
\bkpans{（1）9\quad（2）13\quad（3）6\quad（4）7}

\bkproblem \R{数直線}{すう|ちょく|せん}で、 6と 14は いくつ ちがう？
\bkshikiran
\bkpans{14 － 6 ＝ 8\quad 8}

\bkproblem A、 B、 C の かずを かこう。
\begin{bkfigure}
\begin{tikzpicture}[x=1mm, y=1mm]
  \bkqfline[\footnotesize]{(0,0)}{10.5}{10}{10}{5}
  \bkqfmark{(21,0)}{A}{bkTheme}
  \bkqfmark{(63,0)}{B}{bkTheme}
  \bkqfmark{(94.5,0)}{C}{bkTheme}
\end{tikzpicture}
\end{bkfigure}
\noindent\hfill A\enspace\bkkakikomi[20mm]\hspace{2\zw}B\enspace\bkkakikomi[20mm]\hspace{2\zw}C\enspace\bkkakikomi[20mm]
\bkpans{A\enspace 20\quad B\enspace 60\quad C\enspace 90}

\bkproblem すごろくで、 5の ますに います。 さいころで 6が でると、 どこに すすむ？
\bkshikiran
\bkpans{5 ＋ 6 ＝ 11\quad 11の ます}

\bkproblemgroup{チャレンジ}{1}

\bkproblem 2ばんの えきから 9ばんの えきまで、 なんえき すすむ？
\begin{bkfigure}
\begin{tikzpicture}[x=1mm, y=1mm]
  \draw[bkRoute6, line width=2.6pt, line cap=round] (0,0) -- (110,0);
  \foreach \k in {1,...,10} {\fill[white, draw=bkRoute6!80!black, line width=1.2pt] (\k*11-5.5,0) circle (2.2);
    \node[font=\sffamily\bfseries\footnotesize, text=bkInk] at (\k*11-5.5,-5.5) {\k};}
\end{tikzpicture}
\end{bkfigure}
\bkshikiran
\bkpans{9 － 2 ＝ 7\quad 7えき}

\bkproblem A と B の かずを かこう。
\begin{bkfigure}
\begin{tikzpicture}[x=1mm, y=1mm]
  \bkqfline[\footnotesize]{(0,0)}{10.5}{10}{100}{5}
  \bkqfmark{(31.5,0)}{A}{bkTheme}
  \bkqfmark{(84,0)}{B}{bkTheme}
\end{tikzpicture}
\end{bkfigure}
\noindent\hfill A\enspace\bkkakikomi[24mm]\hspace{3\zw}B\enspace\bkkakikomi[24mm]
\bkpans{A\enspace 300\quad B\enspace 800}
\par\vspace{4pt}

\begin{furikaeri}
\begin{checklist}
\item \R{数直線}{すう|ちょく|せん}で たし\R{算}{ざん}・ひき\R{算}{ざん}が できる。\hfill{\footnotesize\sffamily\color{bkTheme}→\pageref{bk:1-6:susumu}ページ}
\item 2つの かずの ちがいを みつけられる。\hfill{\footnotesize\sffamily\color{bkTheme}→\pageref{bk:1-6:aida}ページ}
\item 1\R{目}{め}もりが 10や 100でも よめる。\hfill{\footnotesize\sffamily\color{bkTheme}→\pageref{bk:1-6:ookii}ページ}
\end{checklist}
チェックが つかなかったら、 →の ページに もどろう。
\end{furikaeri}

\bkfukidashi{sora}{いえの なかで、 \R{数直線}{すう|ちょく|せん}を さがして みよう。 ものさしの ほかにも あるかな？}
\clearpage
